\section{Problem formulation and inferential targets}
\label{sec:problem}

BayesBreak addresses \emph{offline} Bayesian segmentation of ordered data. Throughout, the
fundamental unknown is a contiguous partition of the index set $\{1,\dots,n\}$ together with a
segment parameter attached to each block. This section states the objects of interest in a
model-agnostic way; Section~\ref{sec:notation} introduces the detailed notation used later.

Unless stated otherwise, all inferential statements in the paper use the following standing assumption.

\begin{assumption}[Offline segmentation support and block independence]
\label{ass:standing-offline}
The observations are ordered on a fixed set of candidate boundary locations; the admissible
segment-count set $\{1,\dots,k_{\max}\}$ is finite; every admissible block has finite nonnegative
block evidence, with strictly positive evidence on the support used by the DP and score $0$ (or
log-score $-\infty$) on excluded blocks; segment parameters are conditionally independent across
blocks given a segmentation; and the conditional partition prior factorizes over segments as in
Assumption~\ref{ass:factorizing-prior}. For approximate non-conjugate blocks, the same statement
applies to the approximate score array, while comparison with the exact non-conjugate model
requires the additional error hypothesis of Assumption~\ref{ass:uniform-block-error}.
\end{assumption}

\subsection{Single-sequence segmentation and Bayes regression}
\label{sec:problem-single}

We observe ordered pairs $\{(x_i,y_i)\}_{i=1}^n$ with strictly increasing design points
$x_1<\cdots<x_n$. The design can be irregular. When the observation model includes known
exposures, trial counts, or precision weights, we also allow deterministic nonnegative quantities
$w_i$; unless explicitly stated otherwise, point-observation models use $w_i\equiv 1$.

A segmentation into $k$ contiguous blocks is represented by a boundary vector
\begin{equation}
0=t_0<t_1<\cdots<t_{k-1}<t_k=n,
\end{equation}
where block $q\in\{1,\dots,k\}$ is the index set
\begin{equation}
I_q(t)\coloneqq \{t_{q-1}+1,\dots,t_q\}.
\end{equation}

\begin{definition}[Segmentation space and admissible support]
\label{def:segmentation-space}
For fixed $n$ and $k$, let
\[
\mathcal{T}_{k,n}:=\{t=(t_0,\dots,t_k):0=t_0<t_1<\cdots<t_k=n\}.
\]
The \emph{admissible support} $\mathcal{T}^{+}_{k,n}\subseteq\mathcal{T}_{k,n}$ consists of
those boundary vectors for which every block has positive prior cohesion and finite positive
block evidence. Equivalently, $t\in\mathcal{T}^{+}_{k,n}$ only when
$g(\Delta_x(t_{q-1},t_q))A^{(0)}_{t_{q-1},t_q}>0$ for all $q$, with the obvious replacement of
$A^{(0)}$ by an approximate evidence when a non-conjugate routine is being analyzed as an
approximate score model.
\end{definition}

The corresponding segment parameter is denoted $\theta_q\in\Theta$. Conditionally on
$(k,t,\theta_{1:k})$, observations factorize across blocks:
\begin{equation}
p(y_{1:n}\mid k,t,\theta_{1:k})
= \prod_{q=1}^k \prod_{i\in I_q(t)} p(y_i\mid \theta_q, w_i),
\label{eq:problem-lik}
\end{equation}
where $p(y_i\mid\theta_q,w_i)$ is interpreted according to the chosen family. For example,
depending on the family, $w_i$ encodes exposure in Poisson models, trial count in Binomial models, or known precision
in Gaussian models. Irregular spacing by itself does not force a non-unit $w_i$; design geometry
can instead enter through the partition prior.

For any candidate block $(i,j]$ with $0\le i<j\le n$, the key primitive is the integrated
single-block evidence
\begin{equation}
\mathcal{L}(i,j;\alpha_0,\beta_0)
\coloneqq p(y_{i+1:j}\mid \alpha_0,\beta_0)
= \int \prod_{t=i+1}^j p(y_t\mid\theta,w_t)\,p(\theta\mid\alpha_0,\beta_0)\,d\theta,
\label{eq:problem-block-evidence}
\end{equation}
where $p(\theta\mid\alpha_0,\beta_0)$ is the segment prior with hyperparameters $(\alpha_0,\beta_0)$
(see Section~\ref{sec:notation}). When $\mathcal{L}(i,j)$ is available in closed form, the
entire offline posterior can be computed exactly by dynamic programming. When it is only
approximated, BayesBreak isolates the approximation at the block level and reuses the same global
DP.

\begin{definition}[Boundary-event variables, Bayes curve, and exported estimator]
\label{def:problem-bayes-curve}
For a boundary vector $t\in\mathcal{T}_{k,n}$, define the interior boundary indicator
$b_i(t):=\sum_{p=1}^{k-1}\mathbf{1}\{t_p=i\}$ for $i=1,\dots,n-1$. Because the boundaries are
strictly ordered, $b_i(t)\in\{0,1\}$. For fixed $k$, the Bayes regression curve at index $i$ is
the posterior mean of the observation-scale segment mean covering $i$ after marginalizing over
$t\in\mathcal{T}^{+}_{k,n}$ and over segment parameters. The exported segmentation estimator is a
deterministic boundary vector, usually the joint MAP vector $\widehat t^{\mathrm{MAP}}(k)$, paired
with segment-level posterior summaries for downstream prediction. These two objects need not
coincide because the Bayes curve averages over boundary uncertainty whereas the exported estimator
chooses one segmentation.
\end{definition}

The inferential targets are:
\begin{enumerate}[label=(\roman*), leftmargin=2.3em]
    \item the marginal evidence at fixed segment count,
    \(p(y_{1:n}\mid k)=\sum_{t\in\mathcal{T}^{+}_{k,n}}p(t\mid k)\prod_{q}\mathcal{L}(t_{q-1},t_q)\);
    \item the posterior $p(k,t\mid y)$ over segment counts and boundaries and its marginals
    $p(k\mid y)$, $p(t_p=h\mid y,k)$, and the boundary-event marginal
    \(p(b_i=1\mid y,k)=\sum_{p=1}^{k-1}p(t_p=i\mid y,k)\)
    (the calibration target of \S\ref{sec:experiments});
    \item posterior summaries of the latent piecewise-constant signal, namely
    $p(\theta_i\mid y)$ and $\mathbb{E}[m(\theta_i)\mid y]$ with $m(\theta)$ the
    observation-scale mean parameter (the Bayes regression curve; conditional on the selected
    $\widehat k$ unless stated otherwise);
    \item point estimators for downstream use. In particular the \emph{joint} MAP segmentation
    \(\widehat t^{\mathrm{MAP}}(k)\in\arg\max_t p(t\mid y,k)\)
    differs from the vector of marginal boundary modes
    \(\arg\max_h p(t_p=h\mid y,k)\)
    and is recovered by max-sum DP. The prior on $k$ and the normalizer $\log C_k$ enter the
    objective only when also optimizing over $k$.
\end{enumerate}

The three posterior summaries entering downstream analysis are isolated as labelled objects so
that later sections can refer to them without re-deriving notation.

\begin{definition}[Posterior summaries: marginal boundaries, segment count, and joint MAP]
\label{def:posterior-summaries}
Fix a sequence $y=y_{1:n}$, a maximum segment count $k_{\max}$, and Assumption~\ref{ass:standing-offline}.
Define:
\begin{enumerate}[label=(\alph*), leftmargin=2.3em]
    \item the \emph{marginal boundary posterior}
    \(\mathrm{MB}_i(k):=p(b_i=1\mid y,k)\)
    for interior indices $i\in\{1,\dots,n-1\}$ and admissible $k\in\{1,\dots,k_{\max}\}$;
    \item the \emph{posterior over segment count}
    \(\mathrm{PK}(k):=p(k\mid y)\propto p(k)\,p(y\mid k)\)
    for $k\in\{1,\dots,k_{\max}\}$;
    \item the \emph{joint MAP segmentation at fixed $k$}
    \(\widehat t^{\mathrm{MAP}}(k):=\arg\max_{t\in\mathcal{T}^{+}_{k,n}} p(t\mid y,k)\),
    with ties broken by a deterministic rule (Section~\ref{sec:dp-alg}).
\end{enumerate}
These three objects are returned by distinct DP recursions (sum-product, sum-product with
$k$-aggregation, and max-sum, respectively), and they are not in general consistent with each
other: $\mathrm{MB}_i$ and $\widehat t^{\mathrm{MAP}}(k)$ can place boundary mass at different
indices (Remark~\ref{rem:marg-vs-joint}), and the mode of $\mathrm{PK}$ need not coincide with
the cardinality of any single thresholded set of marginal boundaries.
\end{definition}

\subsection{Multi-sequence structure}
\label{sec:problem-multiseq}

Many applications provide multiple related sequences observed on a common ordered grid or on a
common set of candidate boundary locations. Let sequence $s\in\{1,\dots,S\}$ be denoted by
$y^{(s)}_{1:n}$, with optional weights $w_i^{(s)}$ and subject-specific segment parameters
$\theta_q^{(s)}$ when boundaries are shared. We use $z_s\in\{1,\dots,G\}$ for group membership,
$\mathcal{S}_g:=\{s:z_s=g\}$ for known groups, and $r_{sg}$ for latent-group responsibilities.
BayesBreak treats three progressively richer settings.

\begin{definition}[Multi-sequence segmentation classes]
\label{def:multiseq-classes}
A shared-boundary replicate model has one boundary vector $t$ and subject-specific segment
parameters $\theta^{(s)}_{1:k}$. A known-group model has observed labels $z_s$ and one boundary
vector $t^{(g)}$ per group. A latent-template model has unobserved labels $z_s$ and a finite set of
templates $\tau_g=(k_g,t^{(g)})$ optimized under the template-mixture objective of
Section~\ref{sec:latent-em}. These three classes preserve the block-evidence interface; fully
Bayesian averaging over both unknown labels and all group segmentations is a different target and
is not claimed for the EM template model.
\end{definition}

\paragraph{Shared-boundary replicates.}
All sequences share a common segmentation $t\in\mathcal{T}^{+}_{k,n}$, but each sequence has its
own within-block parameters $\theta^{(s)}_{1:k}$. The common-grid requirement is an assumption on
the candidate boundary coordinates; subject-specific missingness or exposures are handled inside
the block evidences. Conditional on $t$, the sequences are independent:
\begin{equation}
p(\{y^{(s)}\}_{s=1}^S\mid t,\{\theta^{(s)}\})
= \prod_{s=1}^S \prod_{q=1}^k \prod_{i\in I_q(t)} p\bigl(y_i^{(s)}\mid \theta_q^{(s)}, w_i^{(s)}\bigr).
\end{equation}
Integrating out the subject-specific segment parameters yields pooled block evidences obtained by
multiplying subject-level block evidences.

\paragraph{Known groups.}
Sequences are partitioned into known groups $g\in\{1,\dots,G\}$. Each group has its own
segmentation $t^{(g)}\in\mathcal{T}^{+}_{k_g,n}$, while sequences within a group share that
segmentation and retain their own segment parameters. Conditional independence is within group
and across subjects given $t^{(g)}$.

\paragraph{Latent groups.}
Group labels are unknown. A fully Bayesian formulation averages over both
assignments and group-specific segmentations. BayesBreak adopts a tractable
\emph{template-mixture} formulation: each group is associated with a single latent boundary
template $\tau_g=(k_g,t^{(g)})$, and EM alternates between soft assignments (exact responsibilities for the current
templates) and exact responsibility-weighted max-sum DP updates of those templates. The
finite-template objective is monotone non-decreasing under the exact coordinate updates
(Theorem~\ref{thm:em-monotone}); this is a finite-template optimization model rather than exact Bayesian
averaging over latent-group segmentations.

\subsection{Computational scope}
\label{sec:problem-computation}

For fixed $k_{\max}$, the core BayesBreak workflow consists of:
\begin{enumerate}[label=(\alph*), leftmargin=2.3em]
\item precomputing all block evidences and required moment numerators;
\item running sum-product DP recursions for evidences and posterior marginals; and
\item optionally running a max-sum/backtracking recursion for a joint MAP segmentation.
\end{enumerate}
When block evaluation is constant time after prefix-sum preprocessing, this yields
$\mathcal{O}(n^2)$ block precomputation and $\mathcal{O}(k_{\max}n^2)$ DP time. These worst-case
costs make the exact method most natural for moderate sequence lengths (hundreds to low
thousands), where exact posterior quantities are still computationally feasible and often worth
the cost.

\subsection{Prediction as a first-class inferential target}
\label{sec:problem-prediction}

Beyond posterior inference on the training sequence(s), BayesBreak treats \emph{prediction for new
data} as an additional inferential target at the same level as evidence, posteriors, and point
estimators. Given a fitted model $\mathcal{M}_g$ (one per group) and new observations
$(X^{\mathrm{new}},y^{\mathrm{new}})$, the relevant quantities are (i) group-conditional
posterior-predictive likelihoods $p(y^{\mathrm{new}}\mid \mathcal{M}_g)$, (ii) group membership
posteriors $p(g\mid y^{\mathrm{new}})$, and (iii) signal predictions
$\widehat f^{\mathrm{MAP}}_g(x^\star)$ or $\widehat f^{\mathrm{Bayes}}_g(x^\star)$ at query
points $x^\star$. Section~\ref{sec:prediction} develops the prediction layer in detail, covering pointwise,
set-valued, and vector-valued new data, and separates exported-segmentation scoring from optional
full resegmentation of new data.
